\section{A path-local theorem for ordinary flops}
\label{sec:ordinary-flop}

Let \(Y\dashrightarrow Y'\) be a projective ordinary flop covered by
\cite{LeeLinQuWang}, and let \(\ell,\ell'\) be its effective extremal fibre
classes.  The graph correspondence \(F\) identifies the big quantum products
and Poincare pairings after analytic continuation
\[
 q'=q^{-1},\qquad F(\ell)=-\ell'.                         \tag{4.1}
\]
The continuation is analytic in the extremal variable and formal in the
other Novikov variables.

Choose corresponding points \(p\in Y\) and \(p'\in Y'\) outside the
exceptional loci.  Let \(A\) be an ample divisor on the common contraction
and write \(D,D'\) for its pullbacks.  Then
\(D\cdot\ell=D'\cdot\ell'=0\), while \(D\) is positive on every effective
class away from the extremal ray.  We complete the transverse Novikov ring
by \(D\)-degree.

\begin{theorem}[Point-section preservation]
\label{thm:ordinary-flop-point-row}
On each fixed nonsingular continuation domain of \cite{LeeLinQuWang}, the
graph gauge satisfies
\[
 F(s_Y(\OO_p))=s_{Y'}(\OO_{p'}).                         \tag{4.2}
\]
Consequently it preserves the Gamma point row.
\end{theorem}

\begin{proof}
At transverse degree zero, every nonconstant stable map has pure extremal
class and lies in the exceptional locus.  It cannot meet the chosen point.
All positive-extremal descendant invariants with the point insertion
therefore vanish.  The point columns on both sides are their common
classical column for every nonsingular value of \(q\), and the graph
correspondence sends \([p]\) to \([p']\).

Let
\[
 \delta=F(s_Y(\OO_p))-s_{Y'}(\OO_{p'}).
\]
It is jointly flat and has zero transverse constant term.  Suppose
\(\delta_\beta\) is a nonzero coefficient of minimal positive \(D\)-degree.
Horizontality for the Novikov derivation defined by \(D\) gives
\[
 \bigl((D\cdot\beta)\id+z^{-1}(D\star_q-)\bigr)
 \delta_\beta=0.                                           \tag{4.3}
\]
At transverse degree zero, quantum corrections to \(D\star_q\) have pure
extremal degree.  The divisor axiom multiplies them by
\(D\cdot\ell=0\).  Thus \(D\star_q=D\cup-\) in (4.3).  Cup product by \(D\)
is nilpotent and \(D\cdot\beta>0\), so the displayed operator is invertible
over \(\C((z))\).  This contradiction kills the first possible coefficient.
Induction proves (4.2).  Flatness of the pairing then preserves the row
(2.3).
\end{proof}

\begin{corollary}[Packetwise Boolean]
\label{cor:ordinary-flop-packet}
Let \(\cP\subset\mathcal S(Y)\) be a formal-monodromy packet transported by
the Lee--Lin--Qu--Wang graph gauge to
\(\cP'\subset\mathcal S(Y')\).  On the fixed continuation domain,
\[
 b(Y,\cP)=b(Y',\cP').                                      \tag{4.4}
\]
\end{corollary}

\begin{proof}
The graph gauge preserves the quantum connection, grading, first Chern class,
and Poincare pairing.  It therefore transports the complete \(z=0\) formal
type.  Theorem~\ref{thm:ordinary-flop-point-row} identifies the covector on
the transported sectorial realization, and (4.4) follows.
\end{proof}

The theorem is narrower than a Gamma/Fourier--Mukai comparison for every
\(K\)-class.  Pure extremal curves need not miss the support of a general
class, and the support computation in the first paragraph of the proof would
fail.  Nor does the argument assert phase independence outside the chosen
continuation domain.
